\item \textbf{Problema de repaso.} En la operación de un motor de pistón de combustión interna de un solo cilindro, una carga de combustible explota para impulsar el pistón hacia afuera en el llamado tiempo de trabajo. Parte de su energía de salida se almacena en un volante giratorio. Después esta energía se usa para empujar el pistón hacia adentro para comprimir la siguiente carga de combustible y aire. En este proceso de compresión suponga que un volumen original de $0.120\text{ L}$ de un gas ideal diatómico, a presión atmosférica, se comprime adiabáticamente a un octavo de su volumen original. 

(a) Encuentre la entrada de trabajo que se requiere para comprimir el gas. 
(b) Suponga que el volante es un disco sólido de $5.10\text{ kg}$ de masa y $8.50\text{ cm}$ de radio, que gira libremente sin fricción entre el tiempo de trabajo y el tiempo de compresión. ¿Cuán rápido debe girar el volante inmediatamente después del tiempo de trabajo? Esta situación representa la mínima rapidez angular a la que el motor puede funcionar; no es el punto de ahogamiento. 
(c) Cuando la operación del motor está muy arriba del punto de ahogamiento, suponga que el volante pone 5.00\% de su energía máxima para comprimir la siguiente carga de combustible y aire. Encuentre su rapidez angular máxima en este caso.
